% INFO 290 — Group 5 Final Report (DRAFT v5 — compressed for 5-page limit)
% Two-column format. Upload to Overleaf directly.
% Search \TODO and \VERIFY for all pending team action items.
\documentclass[10pt,twocolumn]{article}

\usepackage[margin=0.75in]{geometry}
\usepackage{graphicx}
\usepackage{booktabs}
\usepackage{hyperref}
\usepackage{enumitem}
\usepackage{xcolor}
\usepackage{amsmath}
\usepackage{caption}
\usepackage{subcaption}
\usepackage{natbib}
\usepackage{titlesec}
\usepackage{float}
\usepackage{tabularx}
\usepackage{multirow}
\usepackage{array}

% Spacing for 5-page limit
\setlength{\parskip}{3pt}
\setlength{\parindent}{1em}
\setlength{\columnsep}{0.25in}
\titlespacing*{\section}{0pt}{8pt}{3pt}
\titlespacing*{\subsection}{0pt}{6pt}{2pt}
\titlespacing*{\subsubsection}{0pt}{4pt}{2pt}

% % Draft markers
% \newcommand{\TODO}[1]{\textcolor{red}{\textbf{[TODO: #1]}}}
% \newcommand{\VERIFY}[1]{\textcolor{red}{\textbf{[VERIFY: #1]}}}

\title{\Large Utilization of RAG and Agentic AI\\for Marketing Asset Generation}
\author{
  Ryuichi Ikeda, Priya Charagondla, Shivani Parikh, Nseke Ngilbus, Lena Ehrenreich \\
  University of California, Berkeley \\
  Berkeley, CA 94720, USA \\
  \texttt{\{i.ryuichi,priya.charagondla,neomishivani,nngilbus,lena.ehrenreich\}@berkeley.edu}
}
\date{}

\begin{document}
\maketitle

% ══════════════════════════════════════════════════════════════
\begin{abstract}
We present a system that generates professional marketing 
assets for films—taglines, promotional overviews, and poster 
art direction—by retrieving thematically similar films from a 
corpus of 4,312 TMDB movies and injecting their metadata and 
poster visual descriptions into the generation prompt. We 
compare three system variants: a zero-shot baseline (V1), 
retrieval-augmented generation with visual grounding via 
pre-computed poster captions (V2), and V2 extended with an 
iterative refinement loop using a separate LLM-as-Judge (V3). 
Outputs are evaluated on five dimensions by GPT-4o, using a 
rubric designed around failure modes observed during 
development. Our investigation yields two central findings. 
First, across all three generators tested (Mistral-7B, 
Gemini 2.5 Flash, GPT-4o-mini) and both n=3 and n=20 sample 
sizes, no variant consistently benefits from retrieved 
context: V1 performs comparably to or better than V2 and V3, 
with no statistically significant pairwise differences—a 
result corpus analysis attributes to substantial overlap 
between the textual retrieval surface and content base LLMs 
have likely encountered during training. Second, automated 
evaluation of creative outputs poses unique methodological 
challenges: four of five rubric dimensions score at or near 
ceiling, leaving creative specificity as the only clearly 
discriminating dimension. Rubric tightening partially 
resolves ceiling effects but confirms that low creative 
specificity reflects a genuine generator limitation rather 
than a measurement artifact.
\end{abstract}

% ══════════════════════════════════════════════════════════════
\section{Introduction}

Film marketing assets---taglines, promotional synopses, and poster art direction---are produced through iterative collaboration between creative teams, market analysts, and studio executives. This process is slow, expensive, and relies on tacit knowledge of genre conventions that is difficult to codify. While LLMs can generate creative text, their zero-shot outputs typically lack two properties that professional marketing demands: \textit{strategic grounding} in proven industry patterns, and \textit{genre-specific visual vocabulary} for art direction.

We address this gap by building a system that takes a \textit{structured movie concept} (premise, thematic core, and constraints) and uses \textit{retrieval-augmented generation with visual grounding} to produce a tagline, a promotional overview, and poster art direction. Our approach is guided by two hypotheses: (i)~grounding generation in historical patterns from thematically similar films---including visual descriptions of their poster art---will improve output quality beyond zero-shot generation; (ii)~iterative refinement via an LLM-as-Judge feedback loop can further improve quality by identifying and correcting specific weaknesses. Figure~\ref{fig:journey} illustrates the end-to-end user journey. Section~2 reviews related work; Section~3 describes data sources; Section~4 presents the architecture and variants; Section~5 details experiments; Sections~6 and~7 discuss findings and conclude.

% --- FIGURE 1 PLACEHOLDER ---
\begin{figure}[t]
\centering
\includegraphics[width=\columnwidth]{figures/figure1_user_journey.png}
\caption{End-to-end user journey. A creative professional provides a structured movie concept. The system generates marketing assets, optionally using retrieval (V2/V3) and judge feedback (V3).}
\label{fig:journey}
\end{figure}
% ══════════════════════════════════════════════════════════════
\section{Related Work}

\citet{sun2025ai} compare LLM-generated and human-written marketing slogans, finding comparable consumer engagement but noting that LLM outputs tend toward structural homogeneity and lack market-specific nuance. This identifies a core limitation our work addresses: without external grounding, LLMs produce generic creative text lacking the genre-specific conventions professionals rely on. Our RAG pipeline provides this grounding by retrieving stylistic patterns from historically successful films, and our structured input format with negative constraints gives the generator explicit boundaries to combat homogeneity.

Self-Refine \citep{madaan2023selfrefine} demonstrates that LLMs can iteratively improve their own outputs through self-generated feedback, with diminishing returns beyond 2--3 iterations. Our V3 variant adopts this judge-then-refine architecture, but with a key modification: we separate the judge model (GPT-4o) from the generator (Mistral-7B). Same-model evaluation has been shown to inflate scores \citep{zheng2023judging}; this risk is amplified in creative tasks, where quality is subjective and self-evaluation may confirm the generator's own stylistic tendencies.

% ══════════════════════════════════════════════════════════════
\section{Data Sources and Inputs}
\label{sec:data}

\subsection{TMDB Movie Corpus}
\label{sec:corpus}

We collected metadata for 4{,}816 movies from the TMDB
(The Movie Database) API, including title, tagline, overview,
genre IDs, release date, original language, runtime, audience
metrics (vote average, vote count, popularity), and budget /
revenue figures. The discovery query applied four hard filters
at fetch time: \texttt{vote\_count}\,$\geq$\,100,
\texttt{language}\,=\,\texttt{en-US}, \texttt{status}\,=\,
\texttt{Released}, and a minimum overview length of 50
characters; results were sorted by descending popularity over a
target ceiling of $\sim$5{,}000 titles. After removing 504
records (10.5\%) with missing taglines, the working corpus
contains \textbf{4{,}312 movies}.

The genre distribution is broad but skewed toward Drama (1{,}845)
and Action (1{,}414), followed by Thriller (1{,}366), Comedy
(1{,}221), and Horror (709), with movies multi-labelled at a
mean of 2.7 genres per title. The corpus is 82\% English-original
(3{,}952 titles), with smaller Japanese (181), French (176),
Spanish (93), and Korean (70) representations, and concentrates
in the 2010s and 2020s ($\sim$60\% combined). Each record
contains textual metadata used for retrieval embedding and a
poster image (e.g., \textit{Take Shelter}, 2011, tagline
\textit{``Is it real or just a dream?''}) used for visual
caption generation. The fetch funnel selects for popular,
English-language, post-2000 commercial releases, with
Pearson(\texttt{popularity}, \texttt{vote\_average}) $= 0.098$
indicating that popularity is largely decoupled from quality.
We discuss the implications of this selection bias for RAG
retrieval in \S\ref{sec:discussion}, with a full corpus audit
in Appendix~D.

\subsection{Visual Grounding via Poster Captions}
\label{sec:posters}

Early experiments revealed that without visual references,
generated poster art direction was generic and lacked
compositional detail---e.g., outputs like ``dark and moody
poster'' that no designer could act on. To address this, we
pre-computed structured visual descriptions for all 4{,}312
movie posters. We evaluated three vision-language models across
four configurations (Figure~\ref{fig:vlm}) before selecting
Gemini 2.5 Flash, which produces structured prose across six
visual axes: composition, lighting, color palette, mood,
typography, and symbolic elements. Captions were pre-computed
via the Gemini API prior to any experiments and stored as an
additional column in the retrieval corpus, requiring no API
calls at inference time.

Captions average $\sim$3{,}750 characters---an order of
magnitude longer than the source overviews---and exhibit strong
per-genre distinctive vocabulary (documented in Appendix~D).
Because these AI-generated descriptions are not present in
standard pre-training corpora, the poster-caption channel is
the most likely component of our retrieval surface to supply
non-redundant signal to the generator.

% --- Figure 2: VLM comparison ---
\begin{figure}[t]
\centering
\includegraphics[width=0.30\columnwidth]{figures/shelter.png}

\vspace{0.4em}
\scriptsize
\begin{tabular}{p{1.5cm}p{1.0cm}p{3.4cm}}
\toprule
\textbf{Model} & \textbf{Verdict} & \textbf{Example Output} \\
\midrule
vit-gpt2 & Rejected & ``man standing in front of large group of birds'' \\
BLIP-2 (zero-shot) & Rejected & ``the poster is dark and moody'' \\
BLIP-2 (Q\&A) & Marginal & ``a man and woman in front of a stormy sky'' \\
\textbf{Gemini 2.5 Flash} & \textbf{Adopted} & ``A man with a concerned upward gaze\ldots two funnel clouds on the horizon\ldots'' \\
\bottomrule
\end{tabular}
\caption{VLM captioning comparison on the \textit{Take Shelter} (2011) poster (above). Gemini~2.5~Flash was selected for descriptive richness and accuracy.}
\label{fig:vlm}
\end{figure}

% ══════════════════════════════════════════════════════════════
\section{System Architecture and Variants}

% --- FIGURE 2 PLACEHOLDER ---
\begin{figure}[t]
\centering
\includegraphics[width=\columnwidth]{figures/figure2_system_architecture.png}
\caption{System architecture. V1 generates from input alone. V2 adds retrieved references and poster captions. V3 extends V2 with an iterative judge-feedback loop.}
\label{fig:architecture}
\end{figure}

We implemented three system variants (Figure~\ref{fig:architecture}, Table~\ref{tab:variants}) that share a common output format---a JSON object with tagline, overview, and poster art direction---and differ in how the generation prompt is constructed.

\begin{table}[h]
\centering
\small
\caption{System variant components.}
\label{tab:variants}
\begin{tabular}{lcccc}
\toprule
\textbf{Variant} & \textbf{Input} & \textbf{Retrieval} & \textbf{Gen.} & \textbf{Refine} \\
\midrule
V1 (Zero-Shot) & \checkmark & — & \checkmark & — \\
V2 (RAG+Visual) & \checkmark & \checkmark & \checkmark & — \\
V3 (Agentic) & \checkmark & \checkmark & \checkmark & \checkmark \\
\bottomrule
\end{tabular}
\end{table}

\textbf{Input design (all variants).} Early experiments with free-text prompts produced \textit{narrative contamination}: the generator invented characters, villains, and plot events not specified by the user (e.g., fabricating a disappearance subplot for a premise about a flooded clinic). We replaced free-text input with a three-field structured format: \texttt{core\_premise} (what happens), \texttt{thematic\_core} (the central tension), and \texttt{negative\_constraints} (what the generator must \textit{not} do). With this structure, the generator no longer invents plot elements excluded by user constraints.

\textbf{Retrieval (V2, V3).} Each movie in the corpus is represented by concatenating its title, tagline, and overview, embedding with Sentence-BERT (all-MiniLM-L6-v2, 384~dim), and indexing in FAISS. At query time, the user's premise is embedded and matched against the index, returning the top-$k$ films ($k$ tuned in \S\ref{sec:tuning}). For each retrieved film, a reference block contains its tagline, overview, and pre-computed poster caption. Early experiments showed that retrieval improved stylistic diversity but also introduced \textit{narrative contamination from retrieved plots}---the generator borrowed character names or story events from references. To prevent this, the prompt explicitly instructs the generator to use references for \textit{stylistic and tonal guidance only}.

\textbf{Generation (all variants).} The generator receives the structured input and (for V2/V3) retrieved reference blocks, producing a JSON object with three marketing assets. The prompt instructs the generator to write the overview as compelling promotional copy rather than a plot summary, and to specify concrete visual elements (subject position, lighting, palette, symbolic objects) in the art direction. We use Mistral-7B-Instruct-v0.2 with 4-bit quantization as the base generator. A post-processing step parses and validates the JSON. Implementation details are in Appendix~C.

\textbf{Agentic refinement (V3).} V3 adds an iterative refinement loop. The output is scored by GPT-4o on five dimensions (\S\ref{sec:rubric}). If any dimension scores below~3, the judge identifies the weakest dimension and provides a reason; this feedback is appended to the prompt, and the generator produces a new output. The loop runs for at most 3~iterations, with early stopping when all dimensions pass. The judge (GPT-4o) is separate from the generator (Mistral-7B) to avoid self-evaluation bias \citep{zheng2023judging}.

% ══════════════════════════════════════════════════════════════
\section{Experiments and Evaluation}

\subsection{Evaluation Overview}

Marketing assets are creative artifacts with no single correct answer. Our evaluation must assess whether outputs are \textit{well-crafted}---anchored to the premise, tonally appropriate, visually specific, and creatively distinctive---without comparing to a gold standard. Our evaluation proceeded through rubric design (\S\ref{sec:rubric}), a three-idea case study (\S\ref{sec:casestudy}), generator comparison across three models (\S\ref{sec:generator}), scaled QA-pair generation for statistical evaluation (\S\ref{sec:scaled}), hyperparameter tuning (\S\ref{sec:tuning}), held-out test set evaluation (\S\ref{sec:test}), and ceiling effect analysis (\S\ref{sec:ceiling}).

\subsection{Evaluation Rubric}
\label{sec:rubric}
We defined five evaluation dimensions, each motivated by a specific failure mode observed during development: (1)~\textbf{Narrative fidelity} (motivated by narrative contamination); (2)~\textbf{Genre alignment} (motivated by tone drift); (3)~\textbf{Visual specificity}---concrete, actionable visual elements such as subject position, lighting, palette, symbolic objects (motivated by generic ``dark and moody'' outputs); (4)~\textbf{Creative specificity}---distinctive, non-interchangeable phrasing in the tagline (motivated by the homogeneity problem in \citealp{sun2025ai}); (5)~\textbf{Output format validity}---all JSON fields present and within word limits. Each dimension is scored 1--5 by GPT-4o, with rubric anchors at scores~1, 3, and~5 (full rubric in Appendix~B). The judge is deliberately different from the generator to prevent self-evaluation bias \citep{zheng2023judging}.

\subsection{Case Study ($n{=}3$)}
\label{sec:casestudy}
We began evaluation with three hand-crafted concepts: a Psychological Thriller, a Drama about family and loss, and an Action Thriller. Each was defined with the structured input format and run through all three variants once. These case studies served two purposes: identifying failure modes that informed the rubric, and providing qualitative illustrations of inter-variant differences.

\begin{table}[h]
\centering
\small
\caption{Case study results ($n{=}3$, Mistral-7B). LLM-as-Judge average scores.}
\label{tab:casestudy}
\begin{tabular}{lccc}
\toprule
\textbf{Idea} & \textbf{V1} & \textbf{V2} & \textbf{V3} \\
\midrule
Idea 1 --- Psych.\ Thriller & 4.20 & 4.20 & 4.20 \\
Idea 2 --- Drama & 4.60 & 4.20 & 4.20 \\
Idea 3 --- Action Thriller & 4.60 & 4.20 & 3.80 \\
\bottomrule
\end{tabular}
\end{table}

V1 (zero-shot) performed equal to or better than V2 and V3 in all three cases---unexpected, given that RAG should provide useful grounding. This raised two hypotheses: (a)~the base model (Mistral-7B) may lack the capacity to leverage retrieved context, or (b)~the rubric may not be sensitive enough to detect quality differences. \S\ref{sec:generator} and \S\ref{sec:ceiling} investigate these respectively. Full outputs and posters for Idea~1 are in Appendix~A.

\subsection{Generator Comparison}
\label{sec:generator}
To test hypothesis~(a), we ran the same three ideas through two additional generators (Gemini~2.5~Flash, GPT-4o-mini) and compared the RAG delta (V2$-$V1) across all three models, then extended to a larger validation sample ($n{=}20$).

\begin{table}[h]
\centering
\small
\caption{RAG effect by generator. avg\_$\delta$ = mean(V2$-$V1). Evaluator: GPT-4o (fixed).}
\label{tab:generator}
\begin{tabular}{llccc}
\toprule
\textbf{Generator} & \textbf{n} & \textbf{V1 avg} & \textbf{V2 avg} & \textbf{avg $\delta$} \\
\midrule
Mistral-7B & 3 & 3.80 & 3.47 & $-0.33$ \\
Gemini~2.5~Flash & 3 & 4.73 & 4.93 & $+0.20$ \\
GPT-4o-mini & 3 & 4.20 & 4.20 & $0.00$ \\
\midrule
Gemini~2.5~Flash & 20 & 4.70 & 4.66 & $-0.04$ \\
GPT-4o-mini & 20 & 4.74 & 4.68 & $-0.06$ \\
\bottomrule
\end{tabular}
\end{table}

At $n{=}3$, Gemini~2.5~Flash appeared to benefit from RAG ($+0.20$) while Mistral-7B did not ($-0.33$), consistent with hypothesis~(a). However, at $n{=}20$ the effect disappeared for both API models ($\delta{=}-0.04$ and $-0.06$), indicating the initial Gemini result was noise. Across all three generators and both sample sizes, no variant consistently benefits from retrieved context---the bottleneck is not model capacity alone.

\subsection{Evaluation Data: Scaled QA Pairs}
\label{sec:scaled}
To enable statistical claims, we used GPT-4o to \textit{reverse-generate} structured movie inputs from TMDB metadata: given a real film's title, genre, and overview, generate the structured input (\texttt{core\_premise}, \texttt{thematic\_core}, \texttt{negative\_constraints}) that would produce it. This yielded 200 synthetic inputs. A Mistral-7B quality critic then reviewed each for internal coherence, retaining 181 validated pairs (19 rejected). These were split by stratified sampling on primary genre into validation ($n{=}145$) and held-out test ($n{=}36$) sets. All tuning was conducted on the validation set; the test set was evaluated once with the final configuration.

\subsection{Hyperparameter Tuning}
\label{sec:tuning}
We performed a grid search on the validation set over retrieval count $k \in \{3,5,7\}$, genre filtering $\in \{\text{True, False}\}$, and temperature $\in \{0.4, 0.6, 0.8\}$, yielding 18~configurations evaluated on a random sample of 20~validation records each. The best configuration was $k{=}3$, genre\_filter=False, temperature=0.6 (avg score = 4.50). Notably, genre filtering \textit{reduced} performance, likely because restricting retrieval to same-genre films narrows stylistic diversity.

\subsection{Test Set Evaluation ($n{=}36$)}
\label{sec:test}
Table~\ref{tab:testset} reports results on the held-out test set with the best configuration. All variants use Mistral-7B.

\begin{table}[h]
\centering
\small
\caption{Test set evaluation ($n{=}36$, $k{=}3$, genre\_filter=False, temp=0.6).}
\label{tab:testset}
\begin{tabular}{lcccccc}
\toprule
 & \textbf{Avg} & \textbf{Narr.} & \textbf{Genre} & \textbf{Vis.} & \textbf{Cre.} & \textbf{Fmt.} \\
\midrule
V1 & \textbf{4.79} & 5.00 & 5.00 & 5.00 & \textbf{3.97} & 5.00 \\
V2 & 4.74 & 5.00 & 4.94 & 5.00 & 3.78 & 5.00 \\
V3 & 4.73 & 5.00 & 5.00 & 5.00 & 3.67 & 5.00 \\
\bottomrule
\end{tabular}
\end{table}

V1 achieves the highest overall score (4.79); V2 (4.74) and V3 (4.73) perform comparably. No pairwise difference is statistically significant (V1 vs V2: Wilcoxon signed-rank $p{=}0.06$). Four of five dimensions score at or near ceiling, with creative specificity as the only clearly discriminating dimension (V1$=$3.97, V2$=$3.78, V3$=$3.67). The consistent ordering V1~$>$~V2~$>$~V3 on this dimension suggests that retrieved references and judge feedback do not improve distinctive phrasing with Mistral-7B and may slightly hurt it. V3 fails to outperform V1 for a related reason: if the model cannot effectively utilize retrieved text, natural-language judge feedback places even greater demands on instruction-following.

\subsection{Ceiling Effect Analysis}
\label{sec:ceiling}
Four of five rubric dimensions score near ceiling across all variants: narrative fidelity (5.00), genre alignment (4.94--5.00), visual specificity (5.00), and format validity (5.00). This is both positive---the structured input design (\S4) effectively prevents narrative contamination---and limiting: near-ceiling scores provide no discriminative power between variants, and the overall score differences are driven entirely by creative specificity. The case study (\S\ref{sec:casestudy}), where three qualitatively different outputs received identical scores (4.20), provides concrete evidence of under-discrimination.

We re-judged 20 records with tightened rubric anchors. Genre alignment scores dropped from 5.00 to 4.75--4.95 across all variants, confirming that the original ceiling on this dimension was partly a rubric sensitivity issue. However, creative specificity remained the lowest-scoring dimension across all variants (3.15--3.45) regardless of rubric version, suggesting this reflects a genuine model capability limitation rather than an evaluation artifact. Some ceiling effects are fixable through rubric design; others reflect generator capacity.

% ══════════════════════════════════════════════════════════════
\section{Discussion}
\label{sec:discussion}

\textbf{Largest impact.} The structured input format (Section~4)---not retrieval or refinement---had the largest measurable effect, eliminating narrative contamination and pushing four of five rubric dimensions to ceiling across all variants. This high shared baseline makes inter-variant differences difficult to detect.

\textbf{Most noticeable failure.} V3 degraded creative specificity relative to V1 (3.67 vs.\ 3.97). The refinement loop requires parsing natural-language feedback and selectively revising output---a meta-cognitive demand that Mistral-7B cannot reliably satisfy, causing regression toward generic phrasing.

\textbf{Systematic issues.} Four of five rubric dimensions hit ceiling, concentrating all discriminative signal into creative specificity. Re-judging with tightened anchors confirmed some ceilings are rubric artifacts (genre alignment: 5.00 $\to$ 4.75--4.95) while others reflect genuine generator limitations (creative specificity: 3.15--3.45 regardless of rubric). GPT-4o's dual role as QA generator and judge creates potential evaluation circularity.

\textbf{Robustness and practical trade-offs.} V1's independence from retrieval means it has no failure mode tied to corpus coverage, unlike V2/V3 which could degrade on premises with no thematically similar films in TMDB; robustness to out-of-distribution inputs was not directly tested. Combined with V1's efficiency advantage (\textasciitilde5s vs.\ V3's 20--90s and \textasciitilde\$0.02/run), the zero-shot baseline is the practical choice given the absence of measurable quality benefit from V2/V3.

\textbf{Bias and societal considerations.} Our corpus is 82\% English-original and selection-biased toward popular, post-2000 commercial releases (\S\ref{sec:corpus}, Appendix~D); generated assets reflect these conventions. Selection effects also distort aggregate ratings: non-Anglophone films that survive the popularity filter rate higher than English ones (Japanese 7.40, Korean 7.35 vs.\ English 6.66), simultaneously underrepresenting and over-praising international cinema. Beyond the corpus, Mistral-7B and GPT-4o carry pre-training biases we did not audit---stereotypes or cultural skews in these models propagate silently through the pipeline. Our automated evaluation cannot detect such biases; human review with diverse evaluators would be needed.

% ══════════════════════════════════════════════════════════════
\section{Conclusions and Future Work}
\label{sec:conclusion}

Our key insight is twofold. \textbf{(i)~RAG does not reliably improve over zero-shot for this creative task} across three generators and two sample sizes; retrieved references appear to introduce stylistic noise rather than useful signal. Corpus analysis (Appendix~D) suggests a plausible structural reason: substantial overlap between the textual retrieval surface and content pre-trained LLMs have likely encountered at scale. \textbf{(ii)~Evaluating creative outputs without ground truth requires careful methodology}: rubric design, not just system design, determines what quality differences can be detected---four of five dimensions hit ceiling, and rubric tightening confirmed that some ceilings are measurement artifacts while others (creative specificity) reflect genuine generator limits.

Key limitations: small test set ($n{=}36$), GPT-4o's dual role as input generator and judge, ceiling effects in four of five rubric dimensions, and a corpus selection-biased toward popular English-language releases. With more time and compute, our highest-priority extensions would be (1)~testing a stronger local generator to confirm the model-capacity hypothesis; (2)~human evaluation (blind A/B/C) to validate LLM-as-Judge reliability; (3)~further rubric tightening, especially on creative specificity; and (4)~isolating the poster-caption channel in a RAG ablation, since AI-generated visual descriptions (\S\ref{sec:posters}) are the one retrieval component plausibly outside the base model's pre-training distribution.
% ══════════════════════════════════════════════════════════════
% ══════════════════════════════════════════════════════════════
% ══════════════════════════════════════════════════════════════
\section*{Contributions}

All notebooks are available in our code repository: \url{https://github.com/ryuichi-github/info290-2026s-group5}.

\begin{table}[h]
\centering
\small
\begin{tabularx}{\columnwidth}{l>{\raggedright\arraybackslash}X>{\raggedright\arraybackslash}X}
\toprule
 & \textbf{Sections Written} & \textbf{Code / Notebook} \\
\midrule
Ryuichi & Abstract, \S4 System Architecture and Variants & \textit{data\_pipeline} \\
\midrule
Priya & \S5 Experiments and Evaluation & \textit{experiment\_v2\_completed\_evals} \\
\midrule
Shivani & \S1 Introduction, \S2 Related Work & \textit{agentic\_loop\_movie} \\
\midrule
Nseke & \S6 Discussion, \S7 Conclusions & \textit{nseke\_evaluation (1)}\\
\midrule
Lena & \S3 Data Sources and Inputs & \textit{corpus\_data\_analysis} \\
\bottomrule
\end{tabularx}
\end{table}
%══════════════════════════════════════════════════════════════
\bibliographystyle{plainnat}
\begin{thebibliography}{9}

\bibitem[Sun et~al.(2025)]{sun2025ai}
Sun, L., An, Y., and Song, X. et~al. (2025).
Human vs.\ AI in creative contests: A study on the effectiveness of advertisement copy generated by Chinese large language models.
\textit{Electronic Commerce Research}.

\bibitem[Madaan et~al.(2023)]{madaan2023selfrefine}
Madaan, A., et~al. (2023).
Self-Refine: Iterative Refinement with Self-Feedback.
\textit{NeurIPS 2023}.

\bibitem[Zheng et~al.(2023)]{zheng2023judging}
Zheng, L., et~al. (2023).
Judging LLM-as-a-Judge with MT-Bench and Chatbot Arena.
\textit{NeurIPS 2023}.

\end{thebibliography}

% ══════════════════════════════════════════════════════════════
\clearpage
\onecolumn
\appendix

\section*{Appendix}

\subsection*{A. Case Study: Idea 1 Full Outputs}

\begin{table}[H]
\centering
\small
\caption{Input concept, generated outputs, and poster images for Idea~1 (Psychological Thriller). All variants scored 4.20/5.00 overall by LLM-as-Judge, despite visible qualitative differences.}
\label{tab:idea1_outputs}
\begin{tabularx}{\textwidth}{l>{\raggedright\arraybackslash}X}
\toprule
\textbf{Input} & \textbf{Premise:} A man wakes up in a hospital with no memory of his past. People claiming to be his family and friends appear one by one---each with a different account of who he was. As fragments of memory return, he suspects someone is controlling what he remembers. \newline \textbf{Theme:} Identity is not what you remember---it is what others need you to believe. \newline \textbf{Constraints:} Do not reveal whether visitors are real or fabricated. No action sequences or physical violence. Memory recovery should feel terrifying, not heroic. \\
\bottomrule
\end{tabularx}

\vspace{-0.5em}

\begin{tabularx}{\textwidth}{l>{\raggedright\arraybackslash}X>{\raggedright\arraybackslash}X>{\raggedright\arraybackslash}X}
\toprule
 & \textbf{V1 (Zero-Shot)} & \textbf{V2 (RAG+Visual)} & \textbf{V3 (Agentic)} \\
\midrule
\textbf{Poster} &
\includegraphics[width=\linewidth]{figures/idea1_v1.png} &
\includegraphics[width=\linewidth]{figures/idea1_v2.png} &
\includegraphics[width=\linewidth]{figures/idea1_v3.png} \\
\midrule
\textbf{Tagline} & ``Unraveling the Threads of Identity'' & ``Unraveling Memories'' & ``Who controls your past?'' \\
\midrule
\textbf{Overview} & A man's shattered past resurfaces in a labyrinthine maze of memories, each more deceptive than the last. What if your past was a prison? & Lost in a labyrinth of memories, a man grapples with a shocking revelation: his past may be a twisted tapestry of lies. & A man regains consciousness in a hospital, unsure of his past. Friends and family visit, but their stories contradict one another. \\
\midrule
\textbf{Art Dir.} & Shattered mirror reflecting distorted memories. Endless maze of mirrors. & Expressionless figure surrounded by shattered mirrors reflecting distorted faces. & Man in hospital bed, figures crowd around, faces obscured by shadows. \\
\midrule
\textbf{Score} & 4.20 & 4.20 & 4.20 \\
\bottomrule
\end{tabularx}
\end{table}
\clearpage

\subsection*{B. Full Evaluation Rubric}

\begin{table}[h]
\centering
\small
\caption*{LLM-as-Judge rubric. Scores 2 and 4 are interpolated between anchors. Judge model: GPT-4o (separate from generator).}
\begin{tabularx}{\textwidth}{p{2.3cm}XXX}
\toprule
\textbf{Dimension} & \textbf{Score 5} & \textbf{Score 3} & \textbf{Score 1} \\
\midrule
Narrative Fidelity & No constraints violated. All content strictly anchored to premise and theme. & Minor trope drift; no invented characters or major additions. & Invented characters or plot elements excluded by constraints. \\
\midrule
Genre Alignment & Vocabulary, rhythm, and register fully match genre. & Mostly aligned; one off-genre phrase or image. & Clearly wrong tone. \\
\midrule
Visual Specificity & $\geq$3 of: subject position, lighting, palette, shadow/silhouette, symbolic object. A designer could start sketching. & Atmosphere mentioned but $\geq$2 specifics missing. & Only generic descriptors (``dark and moody''). \\
\midrule
Creative Specificity & Unexpected word combination; memorable on first read. Not interchangeable with other films in this genre. & Competent but interchangeable with many films. & Clich\'{e} phrase from common marketing templates. \\
\midrule
Format Validity & All fields present, within word limits. & One field $\leq$20\% over limit. & Field missing or $>$20\% over. \\
\bottomrule
\end{tabularx}
\end{table}

\subsection*{C. Implementation Details}

\textbf{Generator:} Mistral-7B-Instruct-v0.2, 4-bit NF4 quantization via bitsandbytes, run on Google Colab T4 GPU (16GB VRAM). Sampling parameters: temperature=0.6, top-$p$=0.9, repetition penalty=1.10, max new tokens=320.

\textbf{Retrieval:} FAISS IndexFlatIP (cosine similarity on L2-normalized embeddings). Embedding model: Sentence-BERT all-MiniLM-L6-v2 (384 dimensions).

\textbf{Judge:} GPT-4o via OpenAI API, \texttt{response\_format=json\_object}, max\_tokens=1000.

\textbf{Poster captioning:} Gemini 2.5 Flash via Google AI API. 6-axis prose prompt covering composition, lighting, color palette, mood, typography, and symbolic elements.

\textbf{Prompt template:}
\begin{verbatim}
You are a movie marketing generator specializing in 
distinctive, genre-specific film marketing.

Generate marketing assets for the movie described below.

OUTPUT: ONE JSON object with EXACTLY these keys:
- overview (<= 80 words, compelling promo synopsis)
- tagline (<= 12 words)
- poster_art_direction (<= 60 words)

JSON FORMAT:
{"overview": "", "tagline": "", "poster_art_direction": ""}

MOVIE DETAILS:
Core Premise: __PREMISE__
Thematic Core: __THEME__
IMPORTANT - Do NOT do the following: __NEGATIVES__

REFERENCES from similar successful movies:
STRICT RULE: Use these for WRITING STYLE, TONE, and 
SENTENCE RHYTHM ONLY.
Do NOT borrow any character names, plot events, 
locations, or story elements.
Every word of your output must describe ONLY the movie 
in MOVIE DETAILS above.
If the references do not feel relevant to this film, 
ignore them entirely and generate from the movie 
details alone.
__REFS__

STEP 1 - Study the references above:
- Identify the sentence rhythm and emotional punch 
  of each tagline
- Identify 3 specific visual details from each poster 
  description (lighting, color, composition)
- Note the tone and pacing of each overview
- If references feel stylistically mismatched, 
  skip this step entirely

STEP 2 - Generate assets using ONLY the movie details:
- tagline: Write one unexpected, specific phrase that 
  could ONLY describe this film.
- overview: Open with a hook, not a plot summary. 
  Use present tense.
- poster_art_direction: Specify exactly where the 
  subject is, lighting direction, dominant colors, 
  and one symbolic object.

STRICT RULES:
- Do NOT borrow any character names, plot events, 
  or story elements from references
- Output ONLY the JSON object
- End immediately after the final '}'

Now output the JSON:
\end{verbatim}

\subsection*{D. TMDB Corpus Audit}
\label{app:corpus_audit}

This appendix expands on the corpus characterization summarized
in \S\ref{sec:corpus}--\ref{sec:posters}.

\subsubsection*{D.1 Detailed Distributions}

\textbf{Genre co-occurrence.} Movies are multi-labelled at a
mean of 2.7 genres per title. Pairwise co-occurrence is
concentrated in canonical bundles: Action$+$Thriller (671),
Action$+$Adventure (574), Drama$+$Thriller (560), and
Drama$+$Romance (504), which the retriever can exploit for
stylistic neighbours.

\textbf{Temporal distribution.} Releases concentrate in the
2010s (1{,}470 films) and 2020s (1{,}137); pre-2000 cinema is
comparatively sparse (1{,}148 films across the eight prior
decades). Decadal mean ratings drift downward from 7.64 (1950s)
to 6.62 (2010s), reflecting both rating-inflation for canonized
older films and a markedly larger release volume in the
streaming era; older inputs in the corpus tend to be
\textit{best-of-decade} survivors rather than a representative
cross-section. Tagline coverage declines from 96\% in the
1960s to 83\% in the 2020s.

\textbf{Marketing-surface text statistics.} Tagline texts are
short and stylized: mean 7.4 words and 40 characters, with
4.6\% ending in a question mark and 6.6\% in an exclamation.
Overviews average 268 characters. Financial fields are sparser
by construction: 80.7\% of records report revenue and 76.6\%
report budget, with completeness skewing strongly toward
higher-grossing releases (see D.2).

\textbf{Poster-caption per-genre distinctive vocabulary.} A
per-genre relative-frequency analysis against the corpus
baseline yields highly distinctive vocabularies. Horror posters
foreground \textit{zombie, chucky, insidious, graveyard,
horrifying}; Action skews to \textit{martial, ninja, saiyan,
lightsaber, batman's, indiana}; Thriller to \textit{terminator,
rambo's, wick, lethal, statham's, neeson's}; Crime to
\textit{holmes, heist, handcuffs, ace, illicit}; Sci-Fi to
\textit{godzilla, na'vi, terminator, alien's, extraterrestrial};
Comedy to \textit{comedic, parody, mischief, bewildered,
spongebob}; and Drama to \textit{bittersweet, mentorship,
perseverance, wistful, political}.

\subsubsection*{D.2 Selection Biases in the Fetch Funnel}

The TMDB corpus is not a random sample of cinema; it is the
joint output of a four-stage selection funnel that systematically
over-represents popular, English-language, post-2000 commercial
releases. Each stage contributes a distinct bias with
implications for downstream retrieval.

\textbf{Popularity floor.} The \texttt{vote\_count}\,$\geq$\,100
filter excludes the long tail of niche, festival-circuit, and
cult films---precisely the regime in which retrieval would add
the most marginal information beyond a pre-trained generator's
parametric memory. Within the surviving corpus, audience reach
is itself structured: splitting at the median \texttt{vote\_count}
yields a ``well-known'' half with 96.2\% revenue coverage versus
65.2\% for the ``long-tail'' half, indicating that the financial
fields are themselves a biased subsample.

\textbf{Language and locale.} The English-localised discover
endpoint produces a corpus that is 82\% English-original.
Non-English titles that survive the popularity filter exhibit
higher mean ratings than their English counterparts
(Japanese 7.40, Korean 7.35, Chinese 7.10 vs.\ English
6.66)---a selection effect: only the most successful
international films cross the popularity threshold visible to
an \texttt{en-US} query, so the corpus simultaneously
underrepresents and over-praises non-Anglophone cinema.

\textbf{Temporal skew.} Popularity-sorted discovery oversamples
recent releases (see D.1 above). Older inputs in the corpus
tend to be best-of-decade survivors rather than a representative
cross-section.

\textbf{Quality--popularity decoupling.} The Pearson correlation
between TMDB popularity and \texttt{vote\_average} is
$0.098$---effectively zero. Retrieval ranked by
popularity-correlated similarity therefore does not implicitly
select for higher-quality references; quality and discoverability
operate on largely orthogonal axes.

\subsubsection*{D.3 Implications for RAG}

Taken together, these filters select for content that large
pre-trained language models have likely encountered at scale
during training; TMDB-derived overviews and taglines are common
in public crawls. The retrieval surface for the textual channel
therefore overlaps substantially with the base model's
parametric memory, contributing a structural ceiling on the
marginal information RAG can supply. This is consistent with
the experimental finding on the held-out test set
(\S\ref{sec:test}) that retrieval does not reliably outperform
zero-shot generation in our setup. The poster-caption channel
(\S\ref{sec:posters}) is a partial exception, since its content
is generated rather than scraped---motivating the
poster-caption-only ablation proposed as future work
(\S\ref{sec:conclusion}).

\subsubsection*{D.4 Out-of-Scope Dimensions}

TMDB does not expose voter demographics (age, gender, ethnicity,
country of origin), marketing spend, trailer or social-media
engagement, third-party critical aggregates (e.g., Rotten
Tomatoes, Metacritic), or box-office geography. Claims about
audience segmentation, marketing effectiveness, or critical
reception are therefore beyond what this corpus can support and
would require complementary data sources to investigate.

\end{document}